\section{Jerome's Several Old Testament Projects}

Jerome did not move in one step from the Gospels to a completed Hebrew Old Testament. His Old Testament activity passed through distinguishable Greek-based revisions, sustained study of Hebrew, and a sequence of new translations whose dates and source mixtures differ by book. The common date range about 390--405 describes the main Hebrew-based undertaking; it should not be treated as a single production schedule known in all particulars (Kamesar 1993; Houghton 2023, introduction).

\subsection{Hexaplaric correction before the Hebrew translations}

Jerome had access to the scholarly legacy of Origen's Hexapla: aligned forms of the Hebrew text, Greek transliterations, the Septuagint, and later Jewish Greek versions associated with Aquila, Symmachus, and Theodotion. Letter 32 describes collation work involving Greek and Hebrew materials and revisions of several scriptural books. Its date and vocabulary place it with Greek-based corrective activity, not with a completed translation directly from Hebrew.

Much of Jerome's Hexaplaric revision outside the Psalter did not survive as an independently transmitted whole. That loss prevents a neat reconstruction of every transition. It also warns against reading every later Latin agreement with Greek as an accidental survival: some may descend from deliberate scholarly revision, others from Old Latin ancestry or later correction.

\subsection{Three Psalters, three historical questions}

The Psalms make the layered process unusually visible because Latin worship preserved more than one form.

\begin{fourstable}{Psalter}{Source relation}{Historical role}{Qualification}
Roman & A light Old Latin revision traditionally connected with Jerome & Continued in some liturgical settings, especially at Rome & Jerome's authorship is disputed; ``Roman'' does not prove direct papal production.\\
Gallican & Old Latin revised toward Origen's Hexaplaric Greek text & Became the dominant medieval Vulgate and liturgical Psalter & It is associated with Jerome but is not his direct Hebrew Psalter; early critical signs were not uniformly preserved.\\
\emph{Iuxta Hebraeos} & A new Latin translation based principally on Hebrew, with Greek scholarly aids & Preserved for study and comparison & It did not generally displace the Gallican Psalter in worship.\\
Medieval mixtures & Copying and correction among received forms & Reflect local use and scholarship & A reading requires manuscript-specific analysis; a label alone does not guarantee purity.\\
\end{fourstable}

Jerome's Letter 106, written much later in response to questions about Psalter readings, reveals comparison among Greek and Latin copies, attention to Hebrew, and the afterlife of Hexaplaric critical signs. It is direct evidence of his philological reasoning at that later date and retrospective evidence about the revision's purpose. Michael Graves's edition and commentary and Oliver Norris's survey provide the modern framework for distinguishing the Psalters.

\subsection{The turn \emph{iuxta Hebraeos}}

Around 390 Jerome began translating Old Testament books with Hebrew as the governing source. The undertaking occupied more than a decade. Prefaces, correspondence, and internal comparison permit approximate sequences, but not every traditional book date is equally firm. The Hebrew Psalter belongs early in the movement; Daniel and other books followed in stages; the Pentateuch and some later projects belong toward the end. ``Completed in 405'' is therefore only a convenient boundary for Jerome's principal work, not the date of a single finished Vulgate Bible.

Jerome's Hebrew source was a late-antique consonantal tradition broadly ancestral to the medieval Masoretic Text. It was not identical with the vocalized, accented, and critically edited Hebrew Bibles used today. He also consulted Aquila, Symmachus, Theodotion, Septuagintal and Old Latin forms, exegetical traditions, and Jewish teachers. Adam Kamesar emphasizes the continuing place of Greek scholarship in the Hebrew project. Paul Rodrigue's full case studies of Genesis 37--50, Daniel, and Esther show that Greek Jewish versions, the Septuagint, Old Latin, and Jerome's own interpretation could affect even Hebrew-based translation. Those books establish a mixed source ecology, not one dependency rate for the whole Old Testament. This does not negate Jerome's turn to Hebrew. It specifies what an ancient learned translation from Hebrew entailed.

\subsection{Prefaces as historical instruments}

Jerome's biblical prefaces explain patrons, source choices, book order, textual marks, and objections. They are indispensable, but they are advocacy. A preface may magnify speed, difficulty, novelty, or opposition; a polemical contrast between ``Hebrew truth'' and corrupt ``streams'' advances a case. Modern comparison must test those claims against the translated texts and other witnesses.

The prefaces also traveled with medieval Bibles and taught later readers how to classify books. Their reception could outlast the exact editorial setting that produced them. A statement first written to defend one translation thus became paratext for a composite Bible whose contents and ecclesial canon were not reducible to Jerome's personal judgments.

\claimcheck{``Jerome translated from the Hebrew rather than the Septuagint''}{For much of the Old Testament, Hebrew became his governing source. Yet he had first made Greek-based revisions, continued to consult Greek versions and Latin predecessors, produced more than one Psalter, and did not translate every book in the later Vulgate. The contrast is real but not exhaustive.}
